% Part: sets-functions-relations
% Chapter: relations
% Section: operations

\documentclass[../../../include/open-logic-section]{subfiles}

\begin{document}

\olfileid{sfr}{rel}{ops}
\olsection{संबंधांवरील क्रिया}

संबंधांत बदल करणे किंवा त्यांना एकत्र करणे अनेकदा उपयोगी पडते.
\olref[sfr][rel][ord]{prop:stricttopartial} मध्ये आपण संबंधांचा
\emph{संयोग} पाहिला; तो म्हणजे दोन संबंधांकडे जोड्यांचे संच म्हणून
पाहिल्यावर त्यांचा संयोग. त्याचप्रमाणे,
\olref[sfr][rel][ord]{prop:partialtostrict} मध्ये संबंधांचा सापेक्ष फरक
पाहिला. संबंधांवर करता येणाऱ्या आणखी काही क्रिया पुढीलप्रमाणे आहेत.

\begin{defn}\ollabel{relationoperations}
$R$, $S$ हे संबंध आणि $A$ हा कोणताही संच असू द्या.

$R$ चा \emph{व्युत्क्रम} म्हणजे $R^{-1} = \Setabs{\tuple{y, x}}{\tuple{x,
    y} \in R}$.

$R$ आणि $S$ यांचा \emph{सापेक्ष गुणाकार} म्हणजे $(R \mid S) =
\{\tuple{x, z} : \exists y(Rxy \land Syz)\}$.

$R$ चे $A$ वरील \emph{मर्यादन} म्हणजे $\funrestrictionto{R}{A}= R
\cap A^2$.

$R$ चा $A$ वर \emph{प्रयोग} म्हणजे $\funimage{R}{A} = \{y :
(\exists x \in A)Rxy\}$
\end{defn}

\begin{ex}
$S \subseteq \Int^2$ हा $\Int$ वरील उत्तरवर्ती संबंध असू द्या, म्हणजे
$S = \Setabs{\tuple{x, y} \in \Int^2}{x + 1 = y}$; त्यामुळे $Sxy$
तेव्हा आणि केवळ तेव्हाच, जेव्हा $x + 1 = y$.

$S^{-1}$ हा $\Int$ वरील पूर्ववर्ती संबंध आहे, म्हणजे
$\Setabs{\tuple{x,y}\in\Int^2}{x -1 =y}$.

$S\mid S$ म्हणजे
$ \Setabs{\tuple{x,y}\in\Int^2}{x + 2 =y}$

$\funrestrictionto{S}{\Nat}$ हा $\Nat$ वरील उत्तरवर्ती संबंध आहे.

$\funimage{S}{\{1,2,3\}}$ म्हणजे $\{2, 3, 4\}$.
\end{ex}

\begin{defn}[संक्रमक संवरण] $R \subseteq A^2$ हा द्विपदी संबंध असू द्या.

$R$ चे \emph{संक्रमक संवरण} म्हणजे $R^+ = \bigcup_{0 < n \in
\Nat} R^n$, जेथे $R^1 = R$ आणि $R^{n+1} = R^n \mid R$ अशी पुनरावर्ती
व्याख्या करतो.

$R$ चे \emph{परावर्ती संक्रमक संवरण} म्हणजे $R^* = R^+ \cup \Id{A}$.
\end{defn}

\begin{ex}
$S \subseteq \Int^2$ हा उत्तरवर्ती संबंध घ्या. $S^2xy$ तेव्हा आणि केवळ
तेव्हाच, जेव्हा $x + 2 = y$; $S^3xy$ तेव्हा आणि केवळ तेव्हाच, जेव्हा
$x + 3 = y$; इत्यादी. म्हणून $S^+xy$ तेव्हा आणि केवळ तेव्हाच, जेव्हा
एखाद्या $n \geq 1$ साठी $x + n = y$. दुसऱ्या शब्दांत, $S^+xy$ तेव्हा
आणि केवळ तेव्हाच, जेव्हा $x < y$; आणि $S^*xy$ तेव्हा आणि केवळ तेव्हाच,
जेव्हा $x \le y$.
\end{ex}

\begin{prob}
$R$ चे संक्रमक संवरण खरोखर संक्रमक आहे, हे दाखवा.
\end{prob}

\end{document}
